\chapter{Metodología}

El presente estudio adopta un enfoque metodológico de tipo cuantitativo-experimental, dado que se fundamenta en el análisis de datos objetivos y medibles, tales como la latencia promedio, la latencia p95, el \textit{throughput}, la tasa de error y el consumo de CPU y memoria. Estos datos son recolectados mediante herramientas de monitoreo y observabilidad en entornos controlados.

La unidad de análisis no corresponde a usuarios finales ni a sistemas productivos, sino a microservicios de referencia desplegados y evaluados mediante \textit{$\mu$Bench}, lo que permite simular interacciones y cargas representativas de arquitecturas distribuidas de manera controlada, reproducible y escalable.

El diseño de investigación se estructura en seis fases principales, como se presenta en la Tabla~\ref{tab:fases_metodologicas}.
\begin{table}[htbp]
\centering
\caption{Fases metodológicas del estudio para la evaluación de controles defensivos en microservicios}
\label{tab:fases_metodologicas}

\renewcommand{\arraystretch}{1.3}

\begin{tabular}{p{1cm} p{4cm} p{4.5cm} p{4.5cm}}
\hline
\textbf{Obj.} & \textbf{Fase} & \textbf{Metodología} & \textbf{Herramienta / Recurso} \\
\hline

1 
& Selección de la arquitectura de referencia y controles defensivos 
& Aplicación de un enfoque multicriterio (Weighted Scoring Model y AHP) para la priorización de mecanismos de seguridad relevantes en microservicios 
& Revisión bibliográfica (Rezaei Nasab et al., CSA CCM v4), definición de controles C1--C4 y formalización del diseño experimental factorial \\

2 
& Configuración del entorno experimental 
& Despliegue de microservicios de referencia en Kubernetes, asegurando reproducibilidad y aislamiento experimental 
& Kubernetes (MicroK8s), Prometheus, Grafana y k6 \\

3 
& Implementación de controles defensivos 
& Integración progresiva de mecanismos de seguridad bajo un enfoque DevSecOps, manteniendo comparabilidad entre escenarios 
& Istio o Linkerd (mTLS), Network Policies de Kubernetes, API Gateway (NGINX Ingress, Istio Gateway o Kong), rate limiting \\

4 
& Diseño y ejecución de experimentos 
& Evaluación comparativa entre escenarios base y con controles bajo diseño factorial con bloques aleatorizados (4 controles, 3 variantes por control, 4 niveles de carga y 8 réplicas por celda) 
& k6 (generación de carga), scripts de automatización y matrices de diseño experimental \\

5 
& Recolección y análisis de métricas 
& Análisis cuantitativo de métricas de desempeño y eficiencia (latencia avg/p95, throughput, tasa de error, uso de CPU y memoria). No se incluyen p50 ni p99 porque el objetivo comparativo prioriza tendencia central y cola alta robusta entre tratamientos (avg y p95), evitando percentiles extremos más inestables en ventanas cortas; tampoco se incluye MTTR al no ejecutar en esta fase una campaña formal de inyección de fallos con eventos de recuperación instrumentados. 
& Prometheus (recolección), Grafana (visualización) y Python (procesamiento y análisis) \\

6 
& Interpretación y discusión de resultados 
& Identificación de \textit{trade-offs} entre seguridad, desempeño y eficiencia operativa, contrastando resultados con la literatura 
& Análisis estadístico (ANOVA u otros), visualización de resultados y literatura académica \\

\hline
\end{tabular}

\vspace{0.2cm}
{\raggedright \small \textit{Nota.} La metodología integra selección ex-ante de controles y validación empírica mediante diseño factorial en entornos de microservicios, garantizando trazabilidad, reproducibilidad y consistencia estadística.\par}

\end{table}